\section{Effective action and closure principle}
\label{sec:effective_action}

To move from an auditable lapse proxy to a dynamical spacetime theory, one must specify how protocol observables embed into covariant fields and how those fields enter an effective action.
This step is a \emph{closure}: it cannot be obtained from compilation depth alone, and it is therefore stated explicitly as assumptions in Section~\ref{sec:assumption_ledger}.

\subsection{Action ansatz}
\label{subsec:action_ansatz}

Under Assumptions~\ref{ass:cost_to_density}--\ref{ass:covariant_action} we adopt an effective action of the form
\begin{equation}
  S[g,\chi,\psi_m]
  = \int \dd^4x\,\sqrt{-g}\left[
  \frac{R-2\Lambda}{16\pi G}
  - \lambda_F\, g^{\mu\nu}\nabla_\mu\chi\nabla_\nu\chi
  - V(\chi^2)
  + \mathcal{L}_m(g,\psi_m)
  \right].
  \label{eq:omega_action_cap2}
\end{equation}
Here $R$ is the Ricci scalar, $\Lambda$ is a cosmological term, $\chi$ is the Fisher-amplitude field, and $\mathcal{L}_m$ denotes additional effective matter degrees of freedom.
The potential $V(\chi^2)$ encodes mismatch penalties and model-dependent implementation costs not captured by the quadratic Fisher term.

\begin{remark}[Relation to Fisher information]
\label{rem:fisher_information_relation}
Writing $\rho=\chi^2$, one has
\begin{equation}
  g^{\mu\nu}\nabla_\mu\chi\nabla_\nu\chi
  = \frac{1}{4}\,\frac{g^{\mu\nu}\nabla_\mu\rho\,\nabla_\nu\rho}{\rho}.
\end{equation}
Thus the scalar kinetic term is a covariant Fisher-information penalty on the coarse-grained density profile, a structure familiar from information-theoretic variational principles; see, e.g., \cite{Frieden1998,Reginatto1998}.
\end{remark}

\begin{remark}[Sign choice and ghost-freedom]
\label{rem:lambdaF_sign}
For the scalar sector to be ghost-free in the effective field theory, the kinetic coefficient should satisfy $\lambda_F>0$ (with our mostly-plus convention).
We adopt this sign throughout and treat $\lambda_F$ as a nonnegative fit parameter.
\end{remark}

\begin{remark}[Boundary terms and a well-posed metric variation]
\label{rem:ghy_term}
If the spacetime region has a boundary, the Einstein--Hilbert term requires a boundary contribution for a well-posed Dirichlet variational principle.
One standard choice is the Gibbons--Hawking--York term \cite{York1972,GibbonsHawking1977}
\begin{equation}
  S_{\mathrm{GHY}}
  = \frac{1}{8\pi G}\int_{\partial\mathcal{M}} K\,\sqrt{|h|}\,\dd^3x,
\end{equation}
where $h$ is the induced metric and $K$ is the trace of the extrinsic curvature of the boundary.
In this manuscript we either assume boundary conditions under which the boundary term does not contribute, or we implicitly include $S_{\mathrm{GHY}}$ in the gravitational sector.
\end{remark}

\subsection{Why the metric equation is Einstein}
\label{subsec:einstein_uniqueness}

The left-hand side of the metric field equation is fixed \emph{within the chosen closure class} by the structural requirements of locality, diffeomorphism invariance, and second-order metric equations (Assumption~\ref{ass:covariant_action}).
In four dimensions, Lovelock-type uniqueness implies that the only symmetric, divergence-free rank-2 tensor built from the metric and its derivatives up to second order is $G_{\mu\nu}+\Lambda g_{\mu\nu}$ up to an overall coupling.
Therefore the micro-to-macro content of CAP lies in the \emph{source sector} determined by $\chi$ and $\mathcal{L}_m$, not in modifying the geometric skeleton.
See \cite{Lovelock1971} for the original uniqueness statement and standard GR texts for the effective-field-theory reading.

\begin{remark}[Higher-derivative corrections in an EFT reading]
\label{rem:eft_corrections}
In a generic low-energy effective field theory of gravity, the action includes higher-curvature operators (e.g.\ $R^2$, $R_{\mu\nu}R^{\mu\nu}$) suppressed by a cutoff scale.
CAP-II's \emph{minimal} closure restricts to a second-order metric equation as an explicit modeling choice (Assumption~\ref{ass:covariant_action}); higher-derivative terms can be incorporated systematically as controlled corrections once a readout/cutoff scale is specified.
See, e.g., \cite{Donoghue1994,Burgess2004}.
\end{remark}

\subsection{Regularization, finite parts, and counterterms}
\label{subsec:regularization_counterterms}

Assumption~\ref{ass:canonical_regularization} fixes a canonical rule for extracting finite constants from regulated sums or traces (Abel regularization and finite-part extraction).
At the effective-action level, different finite parts correspond to different choices of local counterterms and hence to different renormalization conditions.
The role of the ledger is to keep this scheme choice explicit and to isolate it from the purely operational definition of $\kappa$ and $\lapse$.


